% Workout Platform --- Bachelor's Thesis Documentation (DRAFT)
% Build: pdflatex main.tex  (run twice for the table of contents / references)
% Only standard packages are used so it builds with a stock TeX Live / MiKTeX.
\documentclass[11pt,a4paper]{article}

\usepackage[utf8]{inputenc}
\usepackage[T1]{fontenc}
\usepackage[margin=2.5cm]{geometry}
\usepackage{booktabs}
\usepackage{enumitem}
\usepackage{xcolor}
\usepackage{listings}
\usepackage{hyperref}
\hypersetup{colorlinks=true, linkcolor=black, urlcolor=blue, citecolor=black}

\definecolor{codebg}{RGB}{245,247,251}
\lstset{
  basicstyle=\ttfamily\footnotesize,
  backgroundcolor=\color{codebg},
  breaklines=true,
  frame=single,
  rulecolor=\color{gray!40},
  columns=fullflexible,
  keepspaces=true,
  showstringspaces=false,
  captionpos=b
}

\title{\textbf{An Evidence-Informed Workout Tracking,\\ Discovery and Coaching Platform}\\[4pt]
\large Design and Implementation of a Full-Stack Training-History System}
\author{Bachelor's Thesis --- Draft Documentation}
\date{\today}

\begin{document}
\maketitle

\begin{abstract}
This document describes the design and implementation of a full-stack workout platform built
as a bachelor's thesis project. The system lets users build evidence-informed workout templates,
execute live training sessions whose history is preserved through immutable snapshots, review
progress through analytics, share and copy templates in a marketplace, and obtain deterministic,
rule-based structural feedback on their programs. The backend is a Spring Boot 3 / Java 21
application backed by PostgreSQL 16 with Flyway migrations; the frontend is a React + TypeScript
single-page application served by the same backend. The work emphasises data integrity
(historical truth via snapshots), correctness (real-database integration tests), security
(JWT authentication with refresh-token rotation and ownership checks), and honest, explainable
feedback rather than opaque ``AI''. This is a working draft and is intentionally incomplete.
\end{abstract}

\tableofcontents
\newpage

\section{Introduction}

\subsection{Motivation}
Most consumer training apps either act as passive loggers or present opaque ``AI coach''
recommendations that users cannot interrogate. The goal of this project is a platform that is
useful as a logger and a planner, that preserves training history truthfully even as plans
change, and that gives \emph{explainable}, evidence-informed structural feedback on programs
without pretending to offer medical or individualised advice.

\subsection{Contributions}
The implemented system provides:
\begin{itemize}[noitemsep]
  \item Authentication with JWT access tokens and rotating refresh tokens.
  \item Management of planning data: a shared official exercise catalog, user custom exercises,
        reusable warm-up/cool-down routines, gyms and equipment.
  \item A multi-day workout \emph{template} builder with snapshotting of exercise metadata.
  \item Live workout execution that copies the chosen template day into \emph{immutable session
        snapshots}, so recorded history never changes when source data is later edited or deleted.
  \item Training history and analytics (volume over time, estimated one-rep-max progression,
        weekly frequency, muscle-group distribution).
  \item A template marketplace (publish, browse, vote, save, deep-copy) built entirely on the
        existing schema.
  \item A deterministic, rule-based, evidence-informed template structure analyzer.
\end{itemize}

\subsection{Scope of this draft}
The platform was built as a sequence of vertical slices (``phases''). This draft documents
phases~1--9 as implemented (Redis caching, phase~8, and OpenSearch/coach mode, phases~10--11,
are out of scope here). It is a work in progress.

\section{Requirements and Scope}

\subsection{Functional requirements}
\begin{itemize}[noitemsep]
  \item Users register, authenticate, and access only their own private data.
  \item Users create and manage exercises, routines, gyms, equipment, and templates.
  \item Users start a workout from a template day at a chosen gym, log sets, add unplanned
        exercises, and finish or cancel the session.
  \item Past sessions remain correct after the source template, exercises, routines, gyms, or
        equipment are edited or deleted.
  \item Users review history and progress analytics derived from real logged data.
  \item Users publish templates publicly, browse and sort the marketplace, vote, save and copy
        templates into their own private library.
  \item Users obtain structural, evidence-informed analysis of a template.
\end{itemize}

\subsection{Non-functional requirements}
PostgreSQL is the single source of truth; the schema is fixed up-front and evolved only through
Flyway migrations. Every user-owned resource enforces ownership checks to avoid Insecure Direct
Object Reference (IDOR) vulnerabilities. The codebase targets ``portfolio / backend systems''
quality rather than a throwaway prototype: no in-memory fakes, no H2, real integration tests.

\subsection{Phased roadmap}
\begin{center}
\begin{tabular}{@{}cll@{}}
\toprule
Phase & Theme & Status \\
\midrule
1 & Infrastructure, database, Docker, health & Done \\
2 & Authentication (JWT + refresh rotation) & Done \\
3 & Planning data (exercises, routines, gyms, equipment) & Done \\
4 & Workout templates / programs & Done \\
5 & Live workout execution + immutable snapshots & Done \\
6 & History + basic analytics & Done \\
7 & Template marketplace & Done \\
8 & Redis caching & Skipped (future) \\
9 & Evidence-informed template analyzer & Done \\
10--11 & Coach mode, OpenSearch & Future \\
\bottomrule
\end{tabular}
\end{center}

\section{System Architecture}

\subsection{High-level shape}
The system is a classic three-tier web application: a React single-page application talks to a
Spring Boot REST API over JSON, and the API persists to PostgreSQL. In production the compiled
frontend is bundled into the backend and served from the same origin at \verb|http://localhost:8080/|,
which keeps the authentication cookie same-origin. The whole stack runs through Docker Compose.

\subsection{Package-by-feature and layering}
The backend is organised \emph{by feature}, not by technical layer. Each feature package
(\verb|auth|, \verb|exercise|, \verb|routine|, \verb|gym|, \verb|template|, \verb|session|,
\verb|history|, \verb|analytics|, \verb|marketplace|, \verb|analyzer|, \verb|shared|) contains its
own four layers:
\begin{itemize}[noitemsep]
  \item \textbf{domain} --- JPA entities and value objects (the persistence model).
  \item \textbf{application} --- services, commands, and feature-specific exceptions (use cases).
  \item \textbf{infrastructure} --- Spring Data repositories.
  \item \textbf{web} --- REST controllers and request/response DTOs.
\end{itemize}
A strict rule is that the web layer never leaks entities: controllers exchange DTOs only. Where
one feature needs another's data (for example the analyzer reading template structure), the
producer exposes a clean \emph{application-level} read model rather than sharing a web DTO or a
repository, keeping the dependency direction sane.

\subsection{Request flow}
A request is authenticated by a servlet filter that validates the JWT and populates the security
context with a \verb|UserPrincipal|. Controllers read the principal, never trusting a user id from
the request body. Services run inside transactions; read paths are marked read-only. Domain errors
extend a shared \verb|ApiException| carrying an HTTP status and a stable machine-readable error
code, and a single \verb|@RestControllerAdvice| renders them into a uniform error envelope.

\subsection{Error envelope}
Every failed request returns a consistent shape:
\begin{lstlisting}
{ "timestamp": "...", "status": 404, "error": "TEMPLATE_NOT_FOUND",
  "message": "Template not found.", "path": "/api/templates/..." }
\end{lstlisting}
The exception handler deliberately omits a catch-all handler so that framework 404s for unknown
routes are not masked; validation failures and unreadable bodies map to \verb|VALIDATION_ERROR|.

\section{Technology Stack}
\begin{center}
\begin{tabular}{@{}ll@{}}
\toprule
Concern & Choice \\
\midrule
Language / runtime & Java 21 \\
Backend framework & Spring Boot 3.5 (Web, Security, Validation, Data JPA) \\
Persistence & PostgreSQL 16, Hibernate ORM 6, \verb|ddl-auto: validate| \\
Migrations & Flyway (\verb|V1| schema, \verb|V2| seed data, \verb|V3| index) \\
Auth & JWT (HS256, jjwt), BCrypt password hashing \\
Frontend & React 18 + TypeScript, Vite, React Router, Recharts \\
Packaging & Docker Compose; frontend bundled into the backend jar \\
Testing & JUnit 5, Spring MVC Test, Testcontainers (real PostgreSQL) \\
\bottomrule
\end{tabular}
\end{center}

\noindent A key discipline is \verb|ddl-auto: validate|: Hibernate never generates or alters
schema. The schema is owned by Flyway, and every entity must match it exactly, which catches
mapping drift at start-up.

\section{Data Model}

\subsection{Overview}
The schema is created by a single initial migration (\verb|V1|) covering users, the exercise
catalog and muscle groups, routines, gyms and equipment, templates and their day/exercise/routine
children, live sessions and their snapshots, the marketplace tables (votes, saves, use events,
per-template stats), and tables reserved for later phases (coaching, analysis results, an outbox).
A second migration seeds the official exercise catalog; a third adds a concurrency index
(Section~\ref{sec:concurrency}).

\subsection{The snapshot principle}
\label{sec:snapshots}
Historical truth is the central data-integrity requirement. It is enforced at two levels:
\begin{enumerate}[noitemsep]
  \item \textbf{Template day exercises} store \verb|exercise_name_snapshot|,
        \verb|exercise_type_snapshot| and a snapshot of muscle groups, keeping only a
        \emph{nullable} reference to the live exercise (\verb|ON DELETE SET NULL|).
  \item \textbf{Workout sessions} snapshot again at start: template name, day name, gym name, and
        every exercise's name/type/planned fields and muscle groups are copied into
        \verb|session_exercises| and \verb|session_exercise_muscle_groups|; logged sets snapshot
        the equipment name.
\end{enumerate}
Because session rows carry their own values and only nullable foreign keys, editing or deleting
the source template, exercise, routine, gym or equipment can never alter what a past session
recorded. This is verified directly by integration tests.

\subsection{Soft deletes and uniqueness}
User-owned planning resources use a \verb|deleted_at| column for soft deletion, so history that
references them survives. Uniqueness is enforced with \emph{partial} unique indexes that ignore
soft-deleted rows, e.g. a unique custom-exercise name per user
\verb|WHERE deleted_at IS NULL|; this lets a name be reused after deletion. Sessions are never
physically deleted --- cancellation is a status, not a row removal.

\section{Security Model}

\subsection{Authentication}
Passwords are stored as BCrypt hashes. On login the API issues a short-lived JWT access token
(sent by the SPA in the \verb|Authorization| header and kept only in memory) and an opaque refresh
token stored as a \verb|HttpOnly|, \verb|SameSite=Strict| cookie scoped to \verb|/api/auth|. The
refresh token is persisted only as a SHA-256 hash. On refresh the token is \emph{rotated}
transactionally: the presented token is validated and revoked and a new pair is issued, so a
stolen-and-replayed refresh token is detected. The cookie's \verb|Secure| flag is profile-based
(off for local HTTP, on for production HTTPS), and CORS allows the Vite dev origin with
credentials.

\subsection{Authorization and IDOR safety}
Every user-owned resource is loaded with an owner-scoped query. Crucially, a request for a resource
owned by another user returns \textbf{404, not 403}, so the API does not even reveal that the
resource exists. Nested resources resolve ownership up the chain (for example a template-day
exercise is reachable only if its day's template is owned by the caller). Public marketplace
templates are readable by any authenticated user; private templates are never exposed in browse
results.

\section{Feature Walkthrough}

\subsection{Planning data}
Users build a library of exercises (the seeded official catalog plus their own custom exercises,
each tagged with primary/secondary muscle groups), reusable START/END routines, gyms, and
equipment within gyms. All endpoints are owner-scoped, validate input, and use soft deletes.

\subsection{Templates}
A template is a private, multi-day program. Each day holds ordered exercises (with planned
sets/reps/weight/rest and a note) and attached routines. Adding an exercise snapshots its
name/type/muscle groups into the template day. After every exercise change the template's
denormalised aggregate arrays (muscle groups, official exercise ids, exercise names) are recomputed
Postgres-side for future search.

\subsection{Live workout execution}
Starting a workout (from a template day and a gym) performs the snapshot copy of
Section~\ref{sec:snapshots} in a single transaction. During the session the user logs sets
(weight, reps, RPE, optional equipment and note), adds unplanned ``extra'' exercises, and finally
finishes or cancels. Only one in-progress workout per user is allowed (Section~\ref{sec:concurrency}).

\subsection{History and analytics}
History lists the user's sessions newest-first with per-session summary numbers. Analytics, computed
entirely with native PostgreSQL aggregation over finished sessions, provide: total volume per day,
weekly workout counts (\verb|date_trunc|), primary-muscle set distribution, and a best estimated
one-rep-max per exercise per day (Epley) for a strength-progression chart.

\subsection{Marketplace}
Owners publish a non-empty template, making it public. Authenticated users browse public templates
sorted by newest, top (net votes) or trending, vote (toggle/switch), save, and \emph{use} a
template. ``Use'' deep-copies the public template into a new private template owned by the copier,
linked to its source only by \verb|copied_from_template_id|. Per-template counters are updated with
atomic, clamped SQL (Section~\ref{sec:concurrency}).

\paragraph{Copying templates with custom exercises.}
A subtle case is a public template that references the publisher's \emph{private} custom exercise.
Because template days are self-contained snapshots, the copy keeps the exercise reference only when
it points to an \emph{official} exercise; for the publisher's custom (or any foreign/deleted)
exercise the reference is set null while the name/type/muscle-group snapshots are preserved. The
copy is therefore fully usable from snapshots and independent of the source, with no privacy leak.

\section{The Evidence-Informed Analyzer}

\subsection{Framing}
The analyzer is deterministic and rule-based --- no machine learning. It produces a \emph{Template
Structure Score}, never a ``scientific quality'' or medical judgement, and every response carries a
disclaimer that it is not medical advice and ignores recovery, technique, effort, sleep, nutrition
and individual response. It is computed live on a read-only \verb|GET| and is not persisted.

\subsection{Inputs and assumptions}
Weekly volume is the sum over the template's authored days (each counted once); when
\verb|days_per_week| exceeds the authored day count an informational note records the assumption.
Sets are weighted by muscle role (primary 1.0, secondary 0.5). Exercises with no planned set count
are \emph{excluded} from volume (never assumed to be a default) and flagged as incomplete data.
Only resistance (\verb|STRENGTH|) exercises count toward volume and balance.

\subsection{Rule families}
The rule pack covers: per-muscle weekly volume bands; major-region coverage (lower body, pulling,
pushing, posterior chain), where minor groups cannot mask a missing major region; weekly frequency
and single-session concentration; structural balance ratios (pull/push, posterior/quads,
lower/upper); rest intervals (compound vs. isolation, the latter detected by a $\geq 3$ distinct
muscle-group heuristic); rep-range extremes (goal-agnostic, since templates have no goal field);
proximity-to-failure (always an informational limitation, as templates store no RIR/RPE); and
difficulty suitability.

\subsection{Scoring}
The 100-point score is split into five sub-scores, each starting at its maximum and deducting a
fixed amount per finding (floored at zero), so outputs are fully deterministic and testable:
\begin{center}
\begin{tabular}{@{}lc l@{}}
\toprule
Sub-score & Max & Example deductions \\
\midrule
Volume coverage & 35 & missing region $-12$; excessive $-5$; very low $-3$ \\
Frequency & 20 & volume concentrated on one day $-4$ \\
Balance & 20 & severe imbalance $-8$; moderate $-4$ \\
Session design & 15 & excessive session $-5$; concentration $-3$ \\
Specificity / rest & 10 & short rest $-2$; rep extreme $-2$ to $-3$ \\
\bottomrule
\end{tabular}
\end{center}
Categories map the total to \verb|WELL_STRUCTURED| ($\geq 80$), \verb|DECENT_STRUCTURE| (55--79) or
\verb|NEEDS_REVIEW| ($<55$).

\section{Engineering Highlights}

\subsection{Concurrency and database constraints}
\label{sec:concurrency}
``One active workout per user'' is guaranteed by a \emph{partial unique index}
(\verb|WHERE status = 'IN_PROGRESS'|) added in migration \verb|V3|; the service-level check is only
a friendly fast path, and a lost race surfaces as a clean \verb|409 ACTIVE_WORKOUT_EXISTS| rather
than a 500. Marketplace vote and save counters are updated with atomic, clamped SQL
(\verb|GREATEST(0, count +/- delta)|) so there is no read--modify--write race and counts never go
negative; counters move only when a row actually changes.

\subsection{Postgres-side aggregation}
Template aggregate arrays and all analytics are computed in SQL rather than in Java. The template
aggregates are recomputed with a single native statement using \verb|array_agg(DISTINCT ... ORDER BY ...)|
with explicitly cast empty-array fallbacks and deterministic ordering (which also keeps tests
stable). This avoids binding PostgreSQL arrays through JPA, which is fragile under
\verb|ddl-auto: validate|.

\subsection{Clean read models}
The analyzer consumes a dedicated application-level read model produced by the template feature,
not a web DTO or a repository, so the analyzer package stays decoupled from template internals.

\section{Testing Strategy}
The backend has 108 integration tests (roughly 3.2k lines) that run against a \emph{real}
PostgreSQL 16 instance via Testcontainers --- never H2 --- so the full Flyway-managed schema and its
constraints are exercised. Tests assert behaviour through the HTTP layer (Spring MVC Test) and, where
it matters, assert database-level guarantees directly: a raw second in-progress insert is rejected
by the partial unique index; recomputed aggregate arrays are read back with \verb|JdbcTemplate|;
session snapshots are shown to be unchanged after the source data is edited. Determinism of the
analyzer is asserted by analysing the same template twice. A representative non-flaky-testing fix
was tracing a JWT-tamper test failure to unused bits in the final base64url character of an HMAC
signature and corrupting a fully-significant character instead.

\section{Limitations and Future Work}
This is an honest list of known gaps:
\begin{itemize}[noitemsep]
  \item \textbf{Frontend tests.} The React frontend (\(\approx\)5.3k lines) currently has no automated
        tests; a component/smoke suite is the most valuable next addition.
  \item \textbf{Auth hardening.} There is no login rate limiting or account lockout.
  \item \textbf{Trending freshness.} \verb|trending_score| is a write-time approximation; a true
        time-decaying rank needs a scheduled recompute (a natural fit for the deferred caching phase).
  \item \textbf{Analyzer validity.} The scoring weights are reasoned heuristics, not empirically
        validated; the feature is framed accordingly and always disclaimed.
  \item \textbf{Search.} Full-text history and marketplace search is reserved for a later
        OpenSearch phase; the denormalised aggregate arrays already prepare for it.
  \item \textbf{Coaching.} Coach/client relationships and coach dashboards are modelled in the
        schema but not yet implemented.
  \item Minor: a small N+1 in the template list, a large frontend bundle (charting library), and
        no generated OpenAPI specification.
\end{itemize}

\section{Conclusion}
The platform realises a coherent, integrity-focused training system: plans are snapshotted into
immutable history, analytics and feedback are derived honestly from real data, and a marketplace
lets programs be shared and independently copied. The engineering emphasises correctness
(real-database tests), data integrity (two-level snapshots, soft deletes), security (rotating JWT
auth and IDOR-safe access), and explainability (a deterministic, disclaimed analyzer) over
superficial features. The remaining phases (caching, coaching, search) extend rather than reshape
this foundation.

\end{document}
